\chapter{Teoria generale della dissonanza}
\label{chap:teoria_dissonanza}

Prima di addentrarci nello studio degli accordi di quattro suoni (le quadriadi), occorre comprendere il comportamento della \textbf{dissonanza} nel sistema tonale. Se la consonanza rappresenta lo stato di quiete e stabilità, la dissonanza è l'elemento dinamico che genera tensione, movimento e la necessità di una risoluzione.

\section{Il ciclo della dissonanza}
\label{sec:ciclo_dissonanza}

Nello stile severo e nella teoria classica dell'armonia, la gestione di una nota dissonante segue un protocollo rigoroso in tre fasi cronologiche.

\begin{regola}[Trattamento della dissonanza]
        \label{reg:trattamento_dissonanza}
        Ogni nota dissonante deve attraversare tre momenti distinti:
        \begin{enumerate}
                \item \textbf{Preparazione:} La nota che diventerà dissonante deve essere presente nell'accordo precedente come \textbf{consonanza} (solitamente mantenuta nella stessa voce). Questa fase "giustifica" l'ingresso della tensione.
                \item \textbf{Percussione:} La nota viene colpita (suonata) simultaneamente alle altre note dell'accordo, diventando effettivamente una \textbf{dissonanza} rispetto alla fondamentale o ad altre voci.
                \item \textbf{Risoluzione:} La nota dissonante deve muoversi per \textbf{grado congiunto discendente} verso una nota che sia consonante nell'accordo successivo.
        \end{enumerate}
\end{regola}

\begin{esempio}[Esempio di preparazione e risoluzione]
        \label{ex:ciclo_dissonanza}
        In un collegamento II-V, la settima del II grado deve essere preparata dalla tonica del I grado o da una nota del IV grado.
        \musica{es_cap08_01}
        In questo esempio, il Do al soprano è consonante nel primo accordo (ottava), diventa dissonante nel secondo (settima del Re) e risolve scendendo al Si (terza del Sol).
\end{esempio}

\section{Evoluzione del concetto di preparazione}
\label{sec:evoluzione_preparazione}

Mentre la \textbf{risoluzione discendente} rimane una regola quasi inviolabile dell'armonia tonale, la necessità della \textbf{preparazione} è andata scemando con l'evoluzione del linguaggio musicale.

Alcune dissonanze, per via della loro forte direzionalità e familiarità (come la settima di dominante), hanno iniziato a essere introdotte "di petto", ovvero senza preparazione. Tuttavia, anche quando la preparazione è omessa, la nota dissonante mantiene intatta la sua natura e la sua esigenza di risoluzione.

\begin{regola}[Dissonanze non preparate]
        \label{reg:dissonanze_non_preparate}
        Gli accordi di "dominante" (\sigla{V}{7}{}, \sigla{VII}{7}{}, e derivati) ammettono frequentemente la percussione della dissonanza senza preparazione. Al contrario, gli accordi di settima sulle altre specie (\sigla{II}{7}{}, \sigla{IV}{7}{}, ecc.) richiedono quasi sempre una preparazione rigorosa nello stile scolastico.
\end{regola}

Questo capitolo è il fondamento teorico di tutti i successivi studi sulle quadriadi:
\begin{itemize}
        \item Per la \textbf{settima di dominante} (Capitolo \ref{chap:settima_dominante}), dove la risoluzione è il motore della cadenza.
        \item Per il \textbf{\sigla{VII}{7}{}} e il \textbf{\sigla{II}{7}{}} (Capitolo \ref{chap:vii7_ii7}), dove la preparazione conta per la fluidità del discorso.
        \item Per tutte le altre \textbf{specie di settima} (Capitolo \ref{chap:sette_specie}), che analizzeremo come estensioni di questo principio.
\end{itemize}
